
\documentclass[11pt]{article}
\usepackage[margin=1in]{geometry}
\usepackage{booktabs}
\usepackage{hyperref}

\title{ RhoeLiquid Rendering Lab Report }
\author{ Example Systems Group }
\date{ 2026-06-01 }

\begin{document}
\maketitle

\section*{Purpose}
Demonstrate that the same Liquid template engine can emit LaTeX artifacts from structured JSON data.

\section*{Protocol}
\begin{enumerate}

  \item Define a small measurement context as JSON.

  \item Render a LaTeX template through the `liquid` CLI.

  \item Compile the generated `.tex` file with a LaTeX toolchain if desired.

\end{enumerate}

\section*{Measurements}
\begin{tabular}{lrrr}
\toprule
Sample & Baseline & Observed & Delta \\
\midrule

alpha & 12.4 & 13.1 & 0.7 \\

beta & 8.8 & 8.4 & -0.4 \\

gamma & 21.0 & 23.5 & 2.5 \\

\bottomrule
\end{tabular}

\section*{Interpretation}

\paragraph{ Signal } All samples rendered through the same template without document-type-specific code.

\paragraph{ Next step } Use a dedicated escaping policy before rendering untrusted content into LaTeX.


\end{document}
